
% !TeX root = ../../../programa/cronograma.tex

% ============================================================
% CRONOGRAMA DE LA EDICIÓN ACTIVA
%
% La primera línea indica a LaTeX Workshop cuál es el documento
% principal que debe compilar cuando este archivo esté abierto.
%
% Este archivo contiene únicamente la tabla específica de una
% edición. La cabecera, bibliografía y configuración general se
% encuentran en programa/cronograma.tex.
% ============================================================


% ------------------------------------------------------------
% CONFIGURACIÓN LOCAL DE LA TABLA
% ------------------------------------------------------------

% \begingroup limita los ajustes siguientes a esta tabla.
\begingroup

% Reduce ligeramente el tamaño del texto.
\small

% Aumenta la altura de las filas.
\renewcommand{\arraystretch}{1.25}

% Controla la separación horizontal entre columnas.
\setlength{\tabcolsep}{5pt}


% ------------------------------------------------------------
% ESTRUCTURA DE LA TABLA
% ------------------------------------------------------------

% longtable permite que la tabla continúe automáticamente
% en páginas sucesivas.
%
% Las cinco columnas corresponden a:
% semana, sesión, tema, lecturas y actividad.
%
% Cada p{...} define un ancho fijo proporcional a \textwidth.
%
% \raggedright alinea a la izquierda.
% \raggedleft alinea a la derecha.
% \arraybackslash permite utilizar \\ al terminar cada celda.
% @{} elimina el espacio automático de los bordes.
\begin{longtable}{@{}
    >{\raggedright\arraybackslash}p{0.055\textwidth}
    >{\raggedright\arraybackslash}p{0.055\textwidth}
    >{\raggedright\arraybackslash}p{0.25\textwidth}
    >{\raggedright\arraybackslash}p{0.44\textwidth}
    >{\raggedleft\arraybackslash}p{0.12\textwidth}
@{}}


% ------------------------------------------------------------
% ENCABEZADO DE LA PRIMERA PÁGINA
% ------------------------------------------------------------

% \toprule crea la línea superior.
\toprule

% & separa columnas y \\ termina la fila.
\textbf{Semana} &
\textbf{Sesión} &
\textbf{Tema} &
\textbf{Lecturas} &
\textbf{Actividad} \\

% Separa el encabezado del contenido.
\midrule

% endfirsthead indica que este encabezado solo aparece en la
% primera página de la tabla.
\endfirsthead


% ------------------------------------------------------------
% ENCABEZADO DE LAS PÁGINAS POSTERIORES
% ------------------------------------------------------------

\toprule

\textbf{Semana} &
\textbf{Sesión} &
\textbf{Tema} &
\textbf{Lecturas} &
\textbf{Actividad} \\

\midrule

% endhead define el encabezado que se repetirá después de cada
% salto de página.
\endhead


% ------------------------------------------------------------
% PIE DE LAS PÁGINAS INTERMEDIAS
% ------------------------------------------------------------

\midrule

% multicolumn combina las cinco columnas en una sola celda.
% La letra r alinea el aviso a la derecha.
\multicolumn{5}{r}{
    \small Continúa en la página siguiente.
} \\

% endfoot define el pie de las páginas intermedias.
\endfoot


% ------------------------------------------------------------
% PIE DE LA ÚLTIMA PÁGINA
% ------------------------------------------------------------

% \bottomrule crea la línea inferior definitiva.
\bottomrule

% endlastfoot define el pie exclusivo de la última página.
\endlastfoot


% ------------------------------------------------------------
% FILA MODELO
% ------------------------------------------------------------

% Duplica y modifica esta estructura para construir el
% cronograma real.
%
% Si una semana contiene dos sesiones, la celda de semana puede
% dejarse vacía en la segunda fila.
%
% \midrule puede utilizarse para separar semanas.
% \addlinespace[4pt] añade espacio entre sesiones sin dibujar
% una línea horizontal.

1 &
I &
Tema de la sesión.
&
Lectura pendiente de definir.
&
Actividad prevista.
\\


% ------------------------------------------------------------
% FIN DE LA TABLA
% ------------------------------------------------------------

\end{longtable}

% Recupera los valores generales anteriores a \begingroup.
\endgroup